\documentclass[letterpaper,11pt]{article}
\usepackage[margin=1.1in]{geometry}

% For angular brackets
\usepackage{textcomp}

% Skip the indents for every paragraph
\usepackage{parskip}

% For customizing itemize spacing
\usepackage{enumitem}
\setlist[itemize]{topsep=0pt, itemsep=0pt, parsep=0pt}
\setlist[enumerate]{topsep=0pt, itemsep=0pt, parsep=0pt}

% For those sweet, sweet acmart fonts
\usepackage{inconsolata}
\RequirePackage[T1]{fontenc} \RequirePackage[tt=false,type1=true]{libertine}

% For more colors!
\usepackage[dvipsnames]{xcolor}

% For URLs
\usepackage[breaklinks,colorlinks]{hyperref}

\begin{document}

\begin{center}
	{\LARGE{\textbf{TA Expectations}}}\\
	  \vspace*{10pt}
	{\Large{CS 408: Software Testing}}\\
        \vspace*{5pt}
	{\Large{Spring 2026}}\\
        \vspace*{5pt}
        Updated: \today \\

  \vspace*{10pt}
\end{center}


% Add regrade policy
% Add group projects different grades

%% Syllabus examples: https://cs17-fall2022.github.io/resources
%%

%% TODO: Check that the regrade policy is here
%% TODO: Say that every group member is responsible for the entire project delivered

\section*{\textcolor{BlueViolet}{Overview}}

Welcome to \underline{CS 408: Software Testing}; I'm looking forward
to working with you this semester! This document lays out the basic
expectations for teaching assistants in this course, so that the
course runs smoothly and the students in the course have consistent
support. As a TA, your duty is to help \emph{facilitate} and
\emph{assess} student learning.

Over the course of the semester, TAs should expect to spend an average
of 10-12 hours per week on course-related duties, with heavier
workloads during weeks when assignments, projects, or exams are being
graded.

\subsection*{Helping Students Learn}
The main avenue for delivering course material is through in-person
lectures, which will be held (mostly) every Tuesday and Thursday from
1:30-2:45 in SMTH 108. The instructor is responsible for
lectures;\footnote{Please let the instructor know if you're
  interesting in giving a guest lecture -- I'm happy for TAs to either
  deliver one of the existing lectures or give a lecture on a topic of
  their choosing!} you're welcome to attend if you want, but are not
required to do so. We will supplement these lectures with an online
discussion board and office hours; which you are expected to help out
with.

\begin{itemize}
\item \textbf{Office Hours:} All members of the course staff will have
  weekly \emph{in-person} office hours. These are optional for
  students and tend to be lightly attended, but they provide an
  important opportunity for in-person assistance from the course
  staff.
  \begin{itemize}
  \item The instructor will have weekly office hours from 12:30-1:30p
    in LWSN 3154L.
  \item Each TA will hold their own (in-person) office hour each week.
  \item During the weeks when students are working on group projects
    (weeks 4-7 and 10-13), each TA is expected to hold an additional
    office hour, for a total of two that week.
  \end{itemize}
  Over the 16 weeks of the semester, each TA will hold a total of 24
  (= 16+ 8) office hours. \emph{Please schedule your office hours by
    the end of the first week of class.} Note that you are free to
  work on whatever you want during this time if no students are
  around!

\item \textbf{Discussion Board:} We will use this
  \href{https://edstem.org/us/courses/91171/discussion}{Ed Discussion}
  site as the course discussion board. In addition to questions about
  course material, students are supposed to post questions about
  projects and homework assignments to the board. \emph{All of the
    course staff, TAs included, is expected to help answer students'
    questions in a timely manner (generally within 24 hours).}  Please
  plan to check the discussion board regularly through the week,
  paying particular attention to questions related to assignments that
  you will be grading. Setting aside a couple of 30-minute blocks over
  the course of the week should be sufficient for this. (Office hours
  are a great time to do this!)
\end{itemize}

\subsection*{Assessing Student Learning}
For better or worse, the bulk of your TA duties will consist of
helping to grade the course's five homework assignments, two group
projects, and two exams. Grading duties will be divided as
follows:
\begin{itemize}
\item Two TAs will each be responsible for \emph{grading one project
    and one homework assignment}. In addition, each of these TAs will
  give a short, in-class introduction to their project the week it is
  released.
\item The remaining TA will be responsible for \emph{grading the other
    three homework assignments}.
\item The entire course staff (including the instructor) will
  collectively grade the midterm and final exams using Gradescope.
\end{itemize}

In order to provide students with timely feedback and an accurate
picture of their standing in the course, we use the following grading
timelines:
\begin{itemize}
\item Homeworks should be graded within a week of their due date,
\item Group projects should be graded within a week and a half of
  their due date, and
\item Exams will be graded within one week
\end{itemize}
\emph{Please think about
  which assignments / projects you would like to be responsible
  for}; our goal is to finalize the grading schedule by the midpoint
of the second week.

\section*{\textcolor{BlueViolet}{Recap}}

Early each week, we will have a brief (<30 minute) meeting to make
sure everything is on track for the upcoming week and to discuss any
issues that have arisen. To summarize, each TA is responsible for:
\begin{itemize}
\item Holding 1-2 in-person office hours every week.
\item Attending the course staff meeting each week.
\item Regularly checking in (2-3x) on the course discussion board
  throughout the week and answering any questions there.
\item Grading either 3 homework assignments or 1 homework assignment
  and 1 group project.
\item Grading a portion of the midterm and final exams.
\end{itemize}
%\newpage
%\bibliographystyle{abbrvnat}
%\bibliography{n-str,p,n}

\end{document}
